\section{The Haim--Kislev Quadratic Program}
\label{sec:quadratic-program-algorithm-hk2019}

The preceding section showed that some capacity-minimizing generalized
Reeb orbit uses a finite list of distinct facet directions. We now express
its action in terms of that list and its dwell times. Haim--Kislev's
formula turns the result into a finite quadratic program~\cite{HK2017}.

The chapter has two tasks. First, it explains the variables and how to
compute the global optimum by solving linear stationarity systems and
comparing their values. Second, it proves reductions for Lagrangian
products of planar polygons. Those reductions make the later family
calculations practical.

\subsection{The finite quadratic program}
\label{subsec:quadratic-program-algorithm-hk2019-general-polytopes}

Let
\[
  K=\{x\in\mathbb{R}^4:\langle a_i,x\rangle\leq 1,\;
    i=1,\ldots,F\}
\]
be a full-dimensional polytope with \(0\in\operatorname{int}(K)\), written as
in Section~\ref{sec:generalized-reeb-orbits-polytopes}. Thus
\(a_i=n_i/h_i\) if \(n_i\) is the outward unit normal to the \(i\)-th facet and
\(h_i=h_K(n_i)>0\) is its height. Haim--Kislev state their formula in the
\((n_i,h_i)\)-variables. We use the equivalent dual-vertex variables
\[
  \beta_i^{\mathrm{here}}=\beta_i^{\mathrm{HK}}h_i.
\]
For a contact-normalized orbit represented by a pure facet word,
$\beta_i$ is the fraction of its period spent
with velocity $2J_0a_i$. The weights therefore sum to one. Closure of the
curve says that the weighted sum of its velocities vanishes; since $J_0$
is invertible, this is equivalent to $\sum_i\beta_i a_i=0$.
Thus the constraints are
\[
  \beta_i\geq0,\qquad
  \sum_i\beta_i=1,\qquad
  \sum_i\beta_i a_i=0.
\]

For an ordered facet word
\(\sigma=(\sigma_1,\ldots,\sigma_m)\) with distinct entries, define
\[
  M_\sigma
  =
  \left\{
    \beta\in\mathbb{R}_{\geq0}^m:
    \sum_{k=1}^m\beta_k=1,\;
    \sum_{k=1}^m\beta_k a_{\sigma_k}=0
  \right\}.
\]
The quadratic objective for this active word is
\begin{equation}\label{eq:quadratic-program-algorithm-hk2019-q}
  Q_\sigma(\beta)
  =
  \sum_{1\leq j<k\leq m}
    \beta_j\beta_k\,\omega_0(a_{\sigma_j},a_{\sigma_k}).
\end{equation}
Here \(\sigma\) is written in the Reeb traversal order used in
Section~\ref{sec:generalized-reeb-orbits-polytopes}. Haim--Kislev's displayed
formula uses the reversed ordered word after translating their symplectic
convention to \(\omega_0\). Since word reversal is a bijection on the ordered
candidate family, the maximum is unchanged; the convention matters only when a
single fixed word is interpreted.

Let
\[
  Q_{\max}
  =
  \max_\sigma\ \max_{\beta\in M_\sigma} Q_\sigma(\beta),
\]
where \(\sigma\) ranges over ordered words with distinct facets and
\(M_\sigma\ne\varnothing\). Equivalently,
one may range over all permutations of all facets and allow zero weights; the
zero-weight entries are then deleted to recover the active word. This is only
an enumeration convention: HK's theorem is the full-permutation statement, and
the active-word form is obtained by deleting unused zero-weight facets and
reindexing the remaining entries. This formulation intentionally permits
duplicates of the same geometric orbit, for example from changing the starting
segment.

\begin{theorem}[Haim--Kislev quadratic program]
\label{thm:quadratic-program-algorithm-hk2019-hk-formula}
For a full-dimensional convex polytope \(K\subset\mathbb{R}^4\) with
\(0\in\operatorname{int}(K)\), written in the facet normalization above,
\(Q_{\max}>0\) and
\[
  c_{\mathrm{EHZ}}(K)=\frac{1}{2Q_{\max}}.
\]
\end{theorem}

\begin{proof}
This is Haim--Kislev's finite capacity formula
\cite[Theorem~1.1]{HK2017}, rewritten from unit normals and heights to the
dual vertices \(a_i=n_i/h_i\). The variable change above sends
\[
  \sum_i\beta_i^{\mathrm{HK}}h_i=1,\qquad
  \sum_i\beta_i^{\mathrm{HK}}n_i=0
\]
to the constraints defining \(M_\sigma\). Their displayed objective uses the
opposite word orientation from \eqref{eq:quadratic-program-algorithm-hk2019-q};
replacing each ordered word by its reverse sends it to \(Q_\sigma\). Since
reversal is a bijection on the words in the maximum, their formula reads
\[
  c_{\mathrm{EHZ}}(K)
  =
  \frac{1}{2}
  \left[\max_\sigma\max_{\beta\in M_\sigma}Q_\sigma(\beta)\right]^{-1}.
\]
\end{proof}

The following construction explains what the finite variables represent and
where global optimality enters.

\Needspace{15\baselineskip}
\begin{proposition}[From finite weights to a dual curve]
\label{prop:quadratic-program-global-maximizer-realization}
Let \(\beta\in M_\sigma\) and put \(q=Q_\sigma(\beta)\).  If \(q>0\), define
\[
  T=\frac{1}{2q},
  \qquad \tau_k=T\beta_k.
\]
Traversing the velocities \(R_{\sigma_k}\) for the durations \(\tau_k\), and
deleting zero-duration pieces, gives a closed simple dual curve \(z\) with
\[
  A(z)=I_K(z)=T.
\]
If \((\sigma,\beta)\) attains \(Q_{\max}\) across all words, then
\(T=c_{\mathrm{EHZ}}(K)\),
and a translation of \(z\) is a minimum-action simple generalized Reeb orbit
on \(\partial K\), with the active word obtained from \(\sigma\) by deleting
its zero-weight entries.

Conversely, if a simple generalized Reeb orbit has period \(T\), active word
\(\sigma\), and positive dwell times \(\tau_k\), then
\(\beta_k=\tau_k/T\) lies in \(M_\sigma\) and
\(Q_\sigma(\beta)=1/(2T)\).
\end{proposition}

\begin{proof}
The closure constraint gives
\[
  \sum_k\tau_kR_{\sigma_k}
  =2T J_0\sum_k\beta_k a_{\sigma_k}
  =0,
\]
while \(\sum_k\tau_k=T\).  The shoelace formula gives
\[
  2A(z)
  =4T^2Q_\sigma(\beta)
  =2T.
\]
Every positive-duration piece has velocity \(R_i\), for which
\[
  H_K^*(-J_0R_i)=\frac14h_K^2(2a_i)=1.
\]
Hence \(I_K(z)=T\), and \(z\) is feasible for the free-period dual problem.

If \((\sigma,\beta)\) attains \(Q_{\max}\), the Haim--Kislev formula gives
\(T=1/(2Q_{\max})=c_{\mathrm{EHZ}}(K)\).  Thus \(z\) is a dual minimizer.
Its multiplier in \eqref{eq:preliminaries-dual-multiplier-value} is
\(\nu=I_K(z)/T=1\).  The reconstruction in
Theorem~\ref{thm:preliminaries-clarke-dual-action} therefore has scale one:
it only translates \(z\) onto \(\partial K\).  Translation preserves the
ordered pure-velocity pieces. On a piece with velocity \(R_i\), the
reconstructed inclusion is
\[
  2a_i=-J_0R_i\in\partial H_K(\gamma)
  =2\operatorname{conv}\{a_j:F_j\text{ is active at }\gamma\}.
\]
The irredundant rows \(a_i\) are extreme points of \(K^\circ\), so this convex
combination forces \(F_i\) to be active there. Thus, after zero durations are
deleted, the resulting minimum-action orbit is simple and has the reduced
active word.

For the converse, closure of the orbit gives
\(\sum_k\beta_k a_{\sigma_k}=0\), and positivity and
\(\sum_k\beta_k=1\) are immediate.  Since \(A(\gamma)=T\), the shoelace
formula yields
\[
  2T
  =
  4T^2
  \sum_{1\leq r<s\leq m}
    \frac{\tau_r}{T}\frac{\tau_s}{T}\,
    \omega_0(a_{\sigma_r},a_{\sigma_s})
  =
  4T^2Q_\sigma(\beta),
  \qquad\text{so}\qquad
  Q_\sigma(\beta)=\frac{1}{2T}.
\]
\end{proof}

\subsection{Solving one word and comparing the results}
\label{subsec:quadratic-program-algorithm-hk2019-computation-layers}

For a fixed word, the constraints are linear and the objective is quadratic.
Write
\[
 C_\sigma=
 \begin{pmatrix}
 a_{\sigma_1}&\cdots&a_{\sigma_m}\\1&\cdots&1
 \end{pmatrix},\qquad e=(0,0,0,0,1)^T,\qquad
 Q_\sigma(\beta)=\tfrac12\beta^T H_\sigma\beta.
\]
Here $H_\sigma$ is symmetric with zero diagonal and, for $j<k$,
$(H_\sigma)_{jk}=(H_\sigma)_{kj}=\omega_0(a_{\sigma_j},a_{\sigma_k})$.
Thus the problem is to maximize this quadratic form subject to
$C_\sigma\beta=e$ and $\beta\geq0$.

First suppose that the maximizing weights are all positive. Small
perturbations in $\ker C_\sigma$ remain feasible, so the first derivative
must vanish on that kernel. Equivalently, there is a vector of multipliers
$\lambda$ such that
\[
 H_\sigma\beta=C_\sigma^T\lambda,\qquad C_\sigma\beta=e.
\]
These are the equality-constrained Karush--Kuhn--Tucker (KKT) equations:
\begin{equation}\label{eq:qp-stationary-system}
 \begin{pmatrix}H_\sigma&-C_\sigma^T\\C_\sigma&0\end{pmatrix}
 \begin{pmatrix}\beta\\\lambda\end{pmatrix}
 =\begin{pmatrix}0\\e\end{pmatrix}.
\end{equation}
They are a linear system even though the objective is quadratic. After
solving it, one checks whether the weights are positive and evaluates $Q$.
A stationary point can be a minimum or a saddle, so this calculation does
not yet identify the maximum.

\begin{example}[A four-facet calculation]
Take $K=[-1,1]^4$ and the ordered facet rows
$(e_1,e_3,-e_1,-e_3)$, where the coordinate order is
$(q_1,q_2,p_1,p_2)$. Closure forces the weights to have the form
$(x,y,x,y)$, and normalization gives $x+y=1/2$. The ordered symplectic
pairings give
\[
 Q=2xy=2x(1/2-x),\qquad 0\leq x\leq1/2.
\]
Thus stationarity reduces to $1-4x=0$. The second derivative is $-4$,
so $x=y=1/4$ maximizes this word, with $Q=1/8$; both boundary choices
have value zero. The corresponding capacity upper bound is $4$.
Identifying the capacity still requires a global
comparison or another argument giving the matching lower bound.
\end{example}

If a maximum has zero weights, delete those entries from the word. The
remaining weights are positive, and deleting zeros changes neither closure
nor the objective. It follows that solving all active words covers maxima
on every face of the nonnegative feasible set. Conversely, a shorter word
can be padded with zero weights to recover a full-facet permutation.
Comparing the feasible stationary values over this complete family
therefore gives the global maximum in Haim--Kislev's formula.

A singular system requires care, but not a search over infinitely many
objective values. Suppose $(\beta,\lambda)$ and
$(\beta+z,\lambda+\mu)$ solve \eqref{eq:qp-stationary-system}. Then
$C_\sigma z=0$ and $H_\sigma z=C_\sigma^T\mu$, whence
\[
 Q_\sigma(\beta+z)-Q_\sigma(\beta)
 =z^TH_\sigma\beta+\tfrac12z^TH_\sigma z
 =\lambda^TC_\sigma z+\tfrac12\mu^TC_\sigma z=0.
\]
Thus all solutions of the same stationary system have the same quadratic
value. One must still decide whether that affine solution set contains
positive weights. A singular matrix by itself means neither infeasibility
nor failure of the capacity formula. This observation will also be used
in the pentagon classification.

A feasible positive value $q$ gives a dual curve with value $1/(2q)$.
Capacity is the infimum of the dual values, so
$c_{\rm EHZ}(K)\leq1/(2q)$. Once $q$ is proved to
be the largest value across all words, that curve is a dual minimizer and
can be translated onto the boundary. For a nonglobal candidate, boundary
placement instead requires the conditions in
Lemma~\ref{lem:generalized-reeb-base-point-recovery}.

The full enumeration can be reduced when a geometric argument proves that
a smaller family still contains a capacity minimizer. For example, an
actual transition from facet $F_i$ to $F_j$ requires an intersection of
those facets and the directed sign $\omega_0(a_i,a_j)\geq0$. These
necessary conditions can exclude words before solving their systems.
They need not hold for every feasible dual curve; what makes them valid
for computing capacity is the existence of a boundary-realized global
minimizer. A stronger reduction for Lagrangian products is proved next.

\subsection{Finite enumeration for Lagrangian products}
\label{subsec:quadratic-program-algorithm-hk2019-lagrangian-products}

Let \(K_q,K_p\subset\mathbb{R}^2\) be nondegenerate convex polygons containing
the origin in their interiors, and let
\[
  K=K_q\times_L K_p\subset\mathbb{R}^2_q\oplus\mathbb{R}^2_p
\]
be their Lagrangian product. A facet of \(K\) is either a \(q\)-facet,
coming from an edge of \(K_q\), or a \(p\)-facet, coming from an edge of
\(K_p\). We call a single facet a one-facet block, and an ordered adjacent pair
of facets of the same type a two-facet block.

\begin{theorem}[Finite enumeration for Lagrangian products]
\label{thm:lagrangian-product-finite-enumeration}
For a Lagrangian product \(K=K_q\times_L K_p\) of nondegenerate planar convex
polygons containing the origin in their interiors, \(c_{\mathrm{EHZ}}(K)\) is
obtained by maximizing \(Q_\sigma\) over the Haim--Kislev quadratic programs
associated to cyclic words formed by $k$ alternating $q$-blocks and
$p$-blocks, where $k\in\{2,3\}$. Each block contains one or two facets
from its factor, distinct blocks use
disjoint facets, and a two-facet block must use adjacent polygon edges. Thus
every word in the family uses at most twelve facets. This is an exhaustive
candidate family for the capacity value; it may contain
overcounting or infeasible words, and it is not a classification of all
minimum-action simple Reeb orbits.
\end{theorem}

\begin{proof}
Put \(T=c_{\mathrm{EHZ}}(K)\). Rudolf's nonsmooth billiard theorem gives a
minimizing Minkowski billiard in \(K_q\), together with its polygonal
dual in \(K_p\), having at most three bounce points and length \(T\)
\cite[Theorem~1]{Rudolf2022}. Write its vertices as \(q_j\) and its dual
vertices as \(p_j\), with indices understood cyclically.  In the notation
needed here, write
\[
  N_C(x)=\{u:\langle u,y-x\rangle\leq0\text{ for every }y\in C\}
\]
for the outward normal cone. Rudolf calls this a strong billiard because
the paired polygons satisfy the following normal-cone relations:
\[
  q_{j+1}-q_j\in N_{K_p}(p_j),
  \qquad
  p_{j+1}-p_j\in-N_{K_q}(q_{j+1}),
\]
Its length is
\[
  \ell_{K_p}(q)
  =\sum_j h_{K_p}(q_{j+1}-q_j)
  =\sum_j\langle p_j,q_{j+1}-q_j\rangle.
\]
The last equality follows from the first normal-cone relation and
\(p_j\in\partial K_p\). These relations say precisely
that the alternating broken line through
\[
  (q_j,p_j),\ (q_{j+1},p_j),\ (q_{j+1},p_{j+1})
\]
is a generalized characteristic for Rudolf's complex-structure convention.
That convention is opposite to ours, so reverse the loop. Since \(\lambda_0\)
differs from \(-p\,dq\) by an exact form, its action in our convention is
\[
  -\oint_{\mathrm{rev}}p\,dq
  =\oint_{\mathrm{orig}}p\,dq
  =\sum_j\langle p_j,q_{j+1}-q_j\rangle
  =\ell_{K_p}(q)=T.
\]
Thus the reversed lift is a capacity-realizing generalized characteristic on
\(\partial(K_q\times_L K_p)\) whose legs alternately move only one factor.
Parametrize it in the contact normalization. Its period is then \(T\), and, as
a dual curve \(z\), it satisfies
\[
  A(z)=I_K(z)=T.
\]

The following surgery is performed on \(z\) as a competitor in the dual
variational problem, not as an orbit already constrained to \(\partial K\).
In particular, splitting a normal-cone direction into pure facet directions
does not assert that the corresponding broken path, drawn directly in
\(K_q\) or \(K_p\), would remain inside that polygon.  Boundary placement is
recovered only after the surgery has been shown to preserve dual minimality.

Apply the splitting--merging--rescaling part of the proof of
Theorem~\ref{thm:generalized-reeb-simple-minimizer} to \(z\). Approximation and
compactness are unnecessary here because the billiard lift is already a finite
broken line. The only extra point to retain from the billiard is its block
bound. On each leg, the polygonal normal cone is generated by one facet normal
or by the two adjacent normals at a vertex, so splitting replaces that leg by
one or two pure directions inside the same block. When a repeated direction is
merged, it is removed from one same-type block and absorbed into another block
that already contains it.  Thus merging cannot increase the number of blocks
of either factor, although it may decrease them.  If the source block is a
singleton, deleting it coalesces its two neighboring opposite-type blocks and
reduces both block counts by one.  This possible loss is harmless: only the
inherited upper bound of three blocks of each type is needed below.

The rescaling argument in the cited proof uses dual minimality to show that the
resulting distinct-facet curve is again a dual minimizer. Indeed, write
\(A'\geq T\) for its action before rescaling; its dual value is still \(T\).
Multiplying every duration by \(T/A'\) gives a dual-feasible curve with
\[
  A=I_K=\frac{T^2}{A'}.
\]
The dual minimum is \(T\), so \(T^2/A'\geq T\), hence \(A'\leq T\). Therefore
\(A'=T\): no rescaling is needed, and the distinct-facet curve \(z'\) is a
dual minimizer with \(A(z')=I_K(z')=T\). By the preceding observations, its
cyclic word still has at most three blocks of each type.

Only now do we return to the primal problem.  Normalize the positive dwell
times of \(z'\) by their sum. The resulting finite weights have QP value
\(1/(2T)=Q_{\max}\), so
Proposition~\ref{prop:quadratic-program-global-maximizer-realization} translates
\(z'\) to a simple generalized Reeb orbit with the same active word on
\(\partial(K_q\times_L K_p)\).  Its factor coordinates therefore remain in
\(K_q\) and \(K_p\); this containment comes from dual-minimizer reconstruction,
not from the splitting operation. A simple orbit cannot contain three
consecutive facets of the same type. Indeed, along
three consecutive \(q\)-facet segments the \(q\)-coordinate is fixed; the two
breakpoints would place it at both endpoints of the middle edge of \(K_q\), a
contradiction. The same holds for \(p\)-facets. Hence every maximal block in
the final word has length one or two, and a two-facet block uses adjacent
edges.

If only one factor type occurred, the whole word would be a single block of
length at most two, whose positive closure equation would require one outward
normal, or a positive combination of two adjacent outward normals, to vanish.
This is impossible. Thus both types occur. Their maximal blocks alternate
cyclically, so there is a common number \(k\leq3\) of \(q\)- and \(p\)-blocks.
One \(q\)-block and one \(p\)-block cannot close, by the same argument in each
factor, so \(k\neq1\). Hence \(k\in\{2,3\}\). Since each block contains at most
two facets, the resulting simple word uses at most \(4k\leq12\) distinct
facets. This proves that the restricted family contains a capacity maximizer.
It does not exclude other minimum-action orbits with more blocks.
\end{proof}

\Needspace{12\baselineskip}
\begin{algorithm}[Lagrangian-product capacity enumeration]
\label{alg:lagrangian-product-capacity-enumeration}
For each factor, enumerate all one-facet blocks and all ordered adjacent
two-facet blocks. For \(k=2,3\), choose \(k\) non-overlapping \(q\)-blocks and
\(k\) non-overlapping \(p\)-blocks, order them cyclically in the alternating
pattern $q$-block, $p$-block, $q$-block, $p$-block, and so on.
Read the facet labels in that order, discard infeasible words, and solve the
corresponding Haim--Kislev quadratic program. The capacity is
\(1/(2Q_{\max})\), where \(Q_{\max}\) is the largest objective value among the
feasible words in this candidate family.
\end{algorithm}

\subsection{A six-facet reduction for the product capacity}
\label{subsec:quadratic-program-product-six-facet}

The product structure gives a reduction that is unavailable for a general
polytope.  The \(q\)- and \(p\)-coordinate planes are Lagrangian, so pairings
between two \(q\)-facets or two \(p\)-facets vanish.  After separating the
total weight in each factor from its normalized distribution, the objective
is therefore linear in either distribution when the other is fixed.
Maximizing these two linear functions at vertices is the mechanism behind
both the six-facet bound and the following algorithm.

Let \(u_1,\ldots,u_r\) be the planar dual vertices of \(K_q\), and let
\(v_1,\ldots,v_s\) be those of \(K_p\).  Define the normalized closure
polytopes
\begin{align*}
  C_q&=\left\{\alpha\in\mathbb R_{\geq0}^{r}:
    \sum_i\alpha_i=1,\ \sum_i\alpha_i u_i=0\right\},\\
  C_p&=\left\{\gamma\in\mathbb R_{\geq0}^{s}:
    \sum_j\gamma_j=1,\ \sum_j\gamma_j v_j=0\right\}.
\end{align*}

\begin{theorem}[Six-facet product maximizer]
\label{thm:quadratic-program-product-six-facet}
For a Lagrangian product \(K_q\times_LK_p\) of nondegenerate planar convex
polygons containing the origin in their interiors, the complete
Haim--Kislev family has a global maximizer using at most three \(q\)-facets
and at most three \(p\)-facets.  The total weight in each factor is \(1/2\).
Consequently the capacity is represented by a word of length at most six.
\end{theorem}

\begin{algorithm}[Product capacity from closure vertices]
\label{alg:quadratic-program-product-closure-vertices}
Enumerate the vertices of \(C_q\) and \(C_p\).  For each pair, enumerate the
cyclic orders of the union of their supports, give each factor total weight
\(1/2\), and evaluate \(Q_\sigma\).  The largest value is \(Q_{\max}\).
Each support union has at most six facets, so fixing one cyclic starting point
leaves at most \(5!=120\) orders.
\end{algorithm}

\begin{figure}[htbp]
  \centering
  \begin{tabular}{c}
    an arbitrary global maximizer \((\sigma,\beta)\)\\[2pt]
    \(\displaystyle\downarrow\quad
      \text{fix \(\gamma\), maximize the linear function in \(\alpha\)}\)\\[2pt]
    \(\alpha\) is a vertex of \(C_q\), hence
      \(\lvert\operatorname{supp}\alpha\rvert\leq3\)\\[2pt]
    \(\displaystyle\downarrow\quad
      \text{fix \(\alpha\), maximize the linear function in \(\gamma\)}\)\\[2pt]
    \(\gamma\) is a vertex of \(C_p\), hence
      \(\lvert\operatorname{supp}\gamma\rvert\leq3\)\\[2pt]
    \(\displaystyle\downarrow\quad
      \text{maximize \(t(1-t)\) and delete zero weights}\)\\[2pt]
    a global maximizer with \(t=\tfrac12\) and at most six facets
  \end{tabular}
  \caption{The proof mechanism for the six-facet reduction.  Linearity in
  each factor distribution is the product-specific step.}
  \label{fig:quadratic-program-product-six-facet-proof}
\end{figure}

\begin{proof}
Choose a global maximizer and let \(I_q\) and \(I_p\) be its \(q\)- and
\(p\)-positions.  Its objective is positive, hence both factor weights are
nonzero.  Put
\[
  t=\sum_{k\in I_q}\beta_k,\qquad
  \alpha_k=\frac{\beta_k}{t}\ (k\in I_q),\qquad
  \gamma_k=\frac{\beta_k}{1-t}\ (k\in I_p).
\]
Closure splits between the factors, so the zero-extended vectors \(\alpha\)
and \(\gamma\) lie in the faces
\[
  C_q(\sigma)=\{\alpha\in C_q:\alpha_i=0
    \text{ if \(q\)-facet \(i\) is omitted from \(\sigma\)}\},
\]
and likewise \(C_p(\sigma)\subset C_p\).  Since the symplectic form vanishes
on each factor, there is a bilinear function
\[
  B_\sigma:C_q(\sigma)\times C_p(\sigma)\longrightarrow\mathbb R
\]
such that
\begin{equation}\label{eq:quadratic-program-product-factorization}
  Q_\sigma(\beta)=t(1-t)B_\sigma(\alpha,\gamma).
\end{equation}
Its coefficients are the signed pairings
\(\langle u_i,v_j\rangle\), with signs determined by the order in \(\sigma\).

With \(\gamma\) fixed, the bilinear factor is linear in \(\alpha\),
so it has a maximizer at a vertex of \(C_q(\sigma)\).  Replacing \(\alpha\)
by that vertex preserves global maximality.  Fixing the new \(\alpha\) and
repeating the argument in \(C_p(\sigma)\) gives a global maximizer for which
both normalized factor weights are vertices of the support-restricted closure
polytopes.  Coordinate restrictions define faces of \(C_q\) and \(C_p\), and
a vertex of a face is a vertex of the ambient polytope.

A vertex of
\[
  \left\{x\geq0:\sum_i x_i=1,\ \sum_i x_iw_i=0\right\},
  \qquad w_i\in\mathbb{R}^2,
\]
has at most three positive coordinates.  Indeed, more than three positive
columns in the three scalar equality constraints have a nontrivial linear
dependence, which gives a feasible line through the point and contradicts
extremality.  Thus each factor uses at most three facets.

For the selected vertices the bilinear value is positive.  Varying only the
total factor weight gives
\[
  \tau(1-\tau)B_\sigma(\alpha,\gamma),\qquad 0\leq\tau\leq1,
\]
which is maximal at \(\tau=1/2\).  Delete the remaining zero-weight positions.
The relative cyclic order and objective are unchanged, and the complete
Haim--Kislev family contains the resulting word of length at most six.

The same argument proves
Algorithm~\ref{alg:quadratic-program-product-closure-vertices}: the selected
vertices occur in the enumeration, while every enumerated candidate is
feasible for the complete Haim--Kislev family and therefore cannot exceed
\(Q_{\max}\).
\end{proof}

\begin{remark}[Why the reduction is product-specific]
The argument does not extend directly to a general polytope.  It uses that
each factor is Lagrangian: same-factor pairings vanish, leaving a bilinear
function which is linear in one normalized closure distribution after the
other is fixed.  Without such a splitting, the objective remains quadratic in
the full weight vector, and moving to a vertex of the closure polytope need not
preserve or improve it.

The theorem asserts the existence of one sparse capacity maximizer.  It neither
classifies all minimizing orbits nor says that every maximizer has support at
most six.  The preceding twelve-facet block family remains useful when an
argument must compare a specified family of geometric branches.  No
genericity assumption is needed: the closure polytopes are closed, and
degeneracy can only lower the rank used in the support bound.
\end{remark}

\subsection{Use in later certificates}
\label{subsec:quadratic-program-algorithm-hk2019-implementation-status}

The later proofs use these results in two ways.
The HKO local-maximum proof uses selected active words to build smooth upper
bounds. The rotated-pentagon proof compares fixed feasible choices at
different angles using the quadratic formula directly.

The numerical implementation follows the same distinction between general
polytopes and products. In the general case it examines the candidate words
and their stationary points, using proved exclusion criteria to reduce the
search. It returns an interval containing the capacity. For a recognized
Lagrangian product it instead enumerates the closure-polytope vertices from
Algorithm~\ref{alg:quadratic-program-product-closure-vertices} and compares
the resulting objective values. This route returns the exact capacity of
the supplied input, whose floating-point coordinates are treated as exact
rational numbers. Both routes use interval bounds to control numerical error
and exact arithmetic when a comparison cannot otherwise be decided.
Section~\ref{sec:numerics} describes those checks. The HKO proof uses a
separate SageMath program over its number field.

The QP is therefore the appropriate language when the global objective value,
its competing branches, or inequalities between them matter.  The next
section keeps the same simple facet words but asks a different geometric
question: it propagates actual breakpoints through facet-pair sections and
detects a closed orbit as a fixed point of an affine return map.  This second
route is conditional on explicit regularity assumptions, but it exposes the
trajectory geometry and provides an independent implementation comparison.
